%======================================================================
%  Hardware platform: SD650 V3 and Intel Xeon CPU Max
%======================================================================

\section{Xeon Max and the SD650 V3}
\label{sec:xeon-max-platform}

\subsection{The unit of analysis}

The Lenovo ThinkSystem SD650~V3 is a dense, direct-water-cooled server
tray rather than a single compute node.  One tray contains two
independent dual-socket nodes, and six trays fit in a 6U DW612S
enclosure~\cite{lp1603}.  This distinction matters when reporting
quantities: the CRESCO8 allocation is a \emph{node}; memory locality is
defined first at the \emph{socket}; and the scheduler may allocate
individual cores within the node.

\begin{table}[tbp]
\centering
\small
\caption{Hardware hierarchy and the properties relevant to this study.
Product capability is separated from the installed CRESCO8
configuration described in Section~\ref{sec:cresco8}
~\cite{lp1603,cresco8doc}.}
\label{tab:hardware-hierarchy}
\begin{tabularx}{\textwidth}{@{}lYY@{}}
\toprule
\textbf{Level} & \textbf{Composition} & \textbf{Relevance} \\
\midrule
DW612S enclosure
  & Six SD650~V3 trays (12 independent nodes)
  & Shared power and cooling infrastructure; not an allocation unit. \\
SD650~V3 tray
  & Two independent dual-socket nodes
  & High-density packaging; the two nodes do not form one shared-memory
    system. \\
Compute node
  & Two Xeon processors connected by Intel UPI
  & Scheduler allocation and the largest cache-coherent NUMA system used
    by the single-node experiments. \\
Xeon Max socket
  & CPU cores, 64\,GB on-package HBM2e and eight DDR5 channels
  & Primary locality domain and the cleanest unit for comparing HBM with
    DDR5. \\
\bottomrule
\end{tabularx}
\end{table}

The product guide covers several processor, memory, storage and network
options.  Most of those options are irrelevant to the thesis.  The
important configuration is the Xeon CPU Max 9480 installed on CRESCO8:
56 physical cores per socket, 64\,GB of HBM2e, eight DDR5-4800 channels,
and four UPI links~\cite{intel9480,lp1603}.  A two-socket node therefore
provides 112 physical cores and 128\,GB of nominal HBM capacity.  The
node also contains DDR5, but its installed capacity and channel
population are properties of CRESCO8, not of the product family.

\begin{figure}[tbp]
\centering
\begin{tikzpicture}[
  font=\small,
  socket/.style={draw=black!65, fill=neutralcol, rounded corners=3pt,
                 minimum width=4.5cm, minimum height=1.35cm,
                 align=center},
  hbm/.style={draw=hbmcol!85!black, fill=hbmcol!18, rounded corners=2pt,
              minimum width=4.5cm, minimum height=0.85cm, align=center},
  ddr/.style={draw=ddrcol!85!black, fill=ddrcol!14, rounded corners=2pt,
              minimum width=4.5cm, minimum height=0.85cm, align=center},
  link/.style={{Stealth}-{Stealth}, thick, black!65}
]
  \node[socket] (s0) at (-3.8,0)
    {\textbf{Socket 0}\\56 cores, AMX, LLC};
  \node[socket] (s1) at (3.8,0)
    {\textbf{Socket 1}\\56 cores, AMX, LLC};
  \node[hbm, below=0.55cm of s0] (h0)
    {64\,GB HBM2e\\high bandwidth};
  \node[hbm, below=0.55cm of s1] (h1)
    {64\,GB HBM2e\\high bandwidth};
  \node[ddr, above=0.55cm of s0] (d0)
    {8 DDR5 channels\\high capacity};
  \node[ddr, above=0.55cm of s1] (d1)
    {8 DDR5 channels\\high capacity};

  \draw[link] (s0) -- node[above, fill=white, inner sep=2pt]{UPI} (s1);
  \draw[link] (s0) -- (h0);
  \draw[link] (s1) -- (h1);
  \draw[link] (s0) -- (d0);
  \draw[link] (s1) -- (d1);

  \node[draw=black!35, dashed, rounded corners=4pt, inner sep=9pt,
        fit=(d0)(d1)(h0)(h1),
        label={[black!60]below:{one cache-coherent compute node}}] {};
\end{tikzpicture}
\caption{Conceptual locality map of a CRESCO8 Xeon Max node.  Each
socket has local HBM and DDR5; access to the other socket traverses UPI.
The diagram represents locality, not physical package routing.
Author's synthesis based on the platform and cluster documentation
~\cite{lp1603,cresco8doc}.}
\label{fig:node-locality}
\end{figure}

\subsection{Compute capability: AMX does not remove the memory problem}

Xeon Max is a Sapphire Rapids processor with conventional x86 cores,
AVX-512, and Intel Advanced Matrix Extensions (AMX).  AMX adds per-core
tile registers and matrix-multiply instructions for BF16 and INT8
operands~\cite{intelmax,na2024llmcpu}.  It can substantially accelerate
matrix operations with suitable shapes and kernels, especially the
large matrix--matrix products encountered during prompt processing.
There is no native INT4 AMX dot-product instruction; low-bit runtimes
must unpack, dequantize, or use specialized integer kernels.  Whether a
model benefits from AMX therefore depends on the runtime, data format,
packing scheme and matrix dimensions.

Peak instruction throughput is not an application result.  Small
matrices, batch-one matrix--vector operations, packing overhead and
memory stalls can leave the matrix units underused.  Likewise, the
advertised maximum turbo frequency is a single-core limit, not the
frequency of a socket executing sustained AMX kernels.  Frequency,
package power and throttling state must be observed during experiments
rather than inferred from specification tables.

\subsection{The memory hierarchy}
\label{sec:memory-hierarchy}

HBM and DDR5 trade different resources.  HBM provides much greater
aggregate bandwidth and 64\,GB of nominal capacity per socket.  DDR5
provides substantially more capacity and can have lower unloaded access
latency.  HBM is therefore not simply ``faster memory'': it is most
valuable when a workload issues enough concurrent requests to approach
its bandwidth and when the bandwidth-critical working set is placed
locally.

Published STREAM results illustrate both the advantage and the need for
local measurement.  Na et al.\ report 588\,GB/s for HBM and
233.8\,GB/s for DDR5 on a Xeon Max 9468 socket
~\cite{na2024llmcpu}.  Ibeid et al.\ report 680--840\,GB/s for HBM and
235--245\,GB/s for DDR5 on a related Xeon Max 9470C
~\cite{ibeid2025aurora}.  These sustained values are 2.5--3.4 times
apart, but the HBM results themselves differ by more than 15\%.  SKU,
clustering mode, thread placement and benchmark configuration are
plausible causes.  None of these numbers should be used as the bandwidth
of CRESCO8 before the Max 9480 nodes are measured.

Three paths must be distinguished in every result:

\begin{itemize}
  \item \textbf{local HBM or local DDR5}, reached from cores on the same
        socket;
  \item \textbf{remote memory}, reached through UPI from the other
        socket; and
  \item \textbf{mixed placement}, in which pages, threads or model
        objects span several locality domains.
\end{itemize}

An apparently node-wide 128\,GB HBM allocation is consequently two
64\,GB local tiers.  A workload uses the aggregate capacity and
bandwidth only if data and computation are balanced across sockets.

\subsection{HBM memory modes}
\label{sec:hbm-modes}

Xeon Max supports three ways to expose HBM
~\cite{lp1603,intelmaxtuning}.  The mode is selected in firmware and is
normally fixed for the lifetime of a boot.

\begin{figure}[tbp]
\centering
\begin{tikzpicture}[
  font=\small,
  cpu/.style={draw=black!65, fill=neutralcol, rounded corners=2pt,
              minimum width=3.9cm, minimum height=0.75cm, align=center},
  hbm/.style={draw=hbmcol!85!black, fill=hbmcol!18, rounded corners=2pt,
              minimum width=3.9cm, minimum height=0.75cm, align=center},
  ddr/.style={draw=ddrcol!85!black, fill=ddrcol!14, rounded corners=2pt,
              minimum width=3.9cm, minimum height=0.75cm, align=center},
  absent/.style={draw=black!35, dashed, text=black!50, rounded corners=2pt,
                 minimum width=3.9cm, minimum height=0.75cm, align=center},
  link/.style={{Stealth}-{Stealth}, thick},
  note/.style={font=\footnotesize, text width=4.2cm, align=center}
]
  \node[font=\bfseries] at (0,1.15) {HBM-only};
  \node[cpu] (c0) at (0,0) {CPU cores};
  \node[hbm] (h0) at (0,-1.25) {HBM};
  \node[absent] (d0) at (0,-2.5) {DDR5 absent};
  \draw[link] (c0) -- (h0);
  \node[note] at (0,-3.55)
    {One visible memory tier; capacity limited to HBM.};

  \node[font=\bfseries] at (5.2,1.15) {Flat};
  \node[cpu] (c1) at (5.2,0) {CPU cores};
  \node[hbm] (h1) at (5.2,-1.25) {HBM NUMA nodes};
  \node[ddr] (d1) at (5.2,-2.5) {DDR5 NUMA nodes};
  \draw[link] (c1) -- (h1);
  \draw[link] (c1.east) .. controls +(0.65,-0.55) and
                                  +(0.65,0.55) .. (d1.east);
  \node[note] at (5.2,-3.55)
    {Both tiers are addressable; software controls placement.};

  \node[font=\bfseries] at (10.4,1.15) {Cache};
  \node[cpu] (c2) at (10.4,0) {CPU cores};
  \node[hbm] (h2) at (10.4,-1.25) {HBM cache};
  \node[ddr] (d2) at (10.4,-2.5) {DDR5 address space};
  \draw[link] (c2) -- (h2);
  \draw[link] (h2) -- (d2);
  \node[note] at (10.4,-3.55)
    {Hardware-managed, direct-mapped cache in front of DDR5.};
\end{tikzpicture}
\caption{The three Xeon Max memory modes, shown for one socket.  Flat
mode is the only mode that exposes HBM and DDR5 as separately
controllable memory tiers.  Author's synthesis from the Xeon Max product
and tuning guides~\cite{lp1603,intelmaxtuning}.}
\label{fig:hbm-modes}
\end{figure}

\begin{description}
  \item[HBM-only mode.] The system contains no DDR5 and exposes HBM as
  ordinary system memory.  It is simple for software but imposes a hard
  capacity limit.  It is not the documented CRESCO8 configuration.

  \item[Flat mode.] HBM and DDR5 share the physical address space but
  appear as distinct NUMA memory nodes.  Applications can allocate
  directly from a tier, or an external policy can bind pages with
  \texttt{mbind}, libnuma or \texttt{numactl}.  Allocation is not enough:
  first touch, page migration and memory mapping can change actual page
  residency, which must be verified.

  \item[Cache mode.] HBM is a hardware-managed, direct-mapped cache in
  front of DDR5.  Applications see the DDR5 capacity without placing
  objects explicitly, but software cannot choose which objects remain
  in HBM.  Conflict misses, warm state and a working set near the cache
  capacity can make performance workload-dependent
  ~\cite{ibeid2025aurora}.
\end{description}

Flat mode provides the strongest experimental design because local HBM
and local DDR5 can be compared on the same socket without changing the
processor or executable.  A Flat-versus-Cache comparison answers a
different question: it measures explicit placement against a hardware
cache policy.

\subsection{Clustering and NUMA}

Memory mode is independent of the processor's clustering mode.  In
Quadrant mode, one socket is normally presented as one compute locality
domain.  Sub-NUMA Clustering (SNC4) exposes four domains per socket,
each associated with a subset of cores, cache, DDR channels and HBM.
SNC4 can increase bandwidth and reduce local latency for software that
places workers and data per quadrant, but it also makes remote accesses
easier to create.  In Flat+SNC4, the nominal 64\,GB HBM socket is divided
into four 16\,GB locality domains.

This distinction explains why a mode cannot be labelled universally
``best''.  Ibeid et al.\ obtain the highest STREAM bandwidth with SNC4
and explicit locality, whereas Na et al.\ observe worse inference
performance and more remote traffic under SNC4
~\cite{ibeid2025aurora,na2024llmcpu}.  The result depends on the
software's placement policy, not only on the processor setting.

\subsection{Cooling and instrumentation}

Direct water cooling allows two high-power processors to operate in a
dense node, but it does not guarantee a fixed frequency or energy
efficiency.  The relevant quantities are measured socket frequency,
temperature, package power and throttling state.  Product-guide claims
about facility heat recovery or cooling savings are not evidence about
inference performance.

The SD650~V3 includes XClarity management and a TI INA226-based
power/energy meter with 100\,Hz sampling capability
~\cite{lp1603}.  This establishes a platform capability, not user
access on CRESCO8.  Intel RAPL, management-controller telemetry and
whole-node power meters cover different system boundaries; they must
not be combined without first validating their scope, units, sampling
and permissions.

\subsection{Consequences for the thesis}

The hardware background yields four requirements for the experimental
work:

\begin{enumerate}
  \item identify the active memory and clustering modes from the
        allocated node rather than from the product guide;
  \item measure sustained local HBM and DDR5 behaviour before using a
        Roofline or interpreting inference results;
  \item bind both threads and pages, then verify page residency and
        remote traffic; and
  \item record frequency, thermal and power state so that a placement
        effect is not confused with a change in operating conditions.
\end{enumerate}
